% ATTEMPT_STATUS: unresolved
% TOP_PROBLEM: {"problemId": "problem.murres-chow-kunneth-conjectures", "canonicalTitle": "Murre's Chow-Kunneth Conjectures", "releaseRank": 153, "primaryDomain": "number_theory_arithmetic_geometry", "primaryDomainLabel": "Number theory & arithmetic geometry", "displayStatus": "open_with_solved_subcases", "releaseStatus": "open_with_solved_subcases"}
% !TeX program = lualatex
% Complete problem section copied verbatim from research_batch_151-153.tex.
\documentclass[11pt]{article}
\usepackage[margin=25mm]{geometry}
\usepackage{fontspec}
\setmainfont{Latin Modern Roman}
\usepackage{amsmath,amssymb,mathtools}
\usepackage{unicode-math}
\setmathfont{Latin Modern Math}
\usepackage{enumitem,needspace}
\usepackage{xurl,hyperref}
\hypersetup{colorlinks=true,urlcolor=blue}
\setcounter{secnumdepth}{1}
\setlength{\parindent}{0pt}
\setlength{\parskip}{5pt}
\setlength{\emergencystretch}{3em}
\setlist[itemize]{leftmargin=3em,itemsep=5pt,topsep=4pt,parsep=0pt}
\allowdisplaybreaks
\newcommand{\N}{\mathbb{N}}
\newcommand{\Z}{\mathbb{Z}}
\newcommand{\Q}{\mathbb{Q}}
\newcommand{\R}{\mathbb{R}}
\newcommand{\C}{\mathbb{C}}
\newcommand{\F}{\mathbb{F}}
\newcommand{\A}{\mathbb{A}}
\newcommand{\PP}{\mathbb P}
\newcommand{\E}{\mathbb{E}}
\newcommand{\eps}{\varepsilon}
\newcommand{\dd}{\,\mathrm d}
\newcommand{\sref}[2]{\hyperlink{source:#1:#2}{[S#2]}}
\newcommand{\eref}[2]{\hyperlink{extra:#1:#2}{[E#2]}}
% Turn the next switch off for a shorter reading copy. The quotations preserve
% every exactTarget field, including qualifications omitted from a short summary.
\newif\ifshowcatalogue
\showcataloguetrue
\newcommand{\cataloguescope}[1]{\ifshowcatalogue
 \paragraph{Exact scope in the supplied catalogue.}
 \begin{quote}\small #1\end{quote}\fi}
\begin{document}
\setcounter{section}{152}
% ================================================================
% problemId: problem.murres-chow-kunneth-conjectures
% source releaseRank: 153; source releaseStatus: open_with_solved_subcases
% exactTarget SHA-256: 981d329e2369d80ecac265c8f104581d62a3a436182899af7698933facc122ac
\clearpage
\section{\texorpdfstring{Murre's Chow-Kunneth Conjectures}{Murre's Chow-Kunneth Conjectures}}\label{153}
\subsection{Definitions and mathematical statement}
For a smooth projective complex $d$-fold $X$, $\mathrm{CH}^j(X)_{\Q}$ is the rational vector space of codimension-$j$ cycles modulo rational equivalence. Degree-zero correspondences belong to $\mathrm{CH}^d(X\times X)_{\Q}$ and act by pullback, intersection, and pushforward. A Chow--K\"unneth decomposition consists of orthogonal idempotents $\pi_i$, summing to the diagonal, whose cohomological actions project onto $H^i(X,\Q)$. Murre's package asks for such a decomposition, vanishing of $(\pi_i)_*$ on $\mathrm{CH}^j$ outside $j\le i\le2j$, and a decomposition-independent filtration
\[
 F^r\mathrm{CH}^j=\bigcap_{i=\max(0,2j-r+1)}^{2j}\ker(\pi_i)_*,
 \qquad F^1\mathrm{CH}^j=\mathrm{CH}^j_{\mathrm{hom}}.
\]
An empty intersection is the whole group; the homologically trivial subgroup is the kernel of the Betti cycle-class map. The universal quantification over \emph{every} decomposition in B--D is retained, not replaced by existence of one specially chosen decomposition.
\par\smallskip\textit{Formulation and concept references:} \sref{153}{1}.
\cataloguescope{For every smooth projective irreducible complex variety X of dimension d, let CH\textasciicircum{}j(X)\_Q denote cycles modulo rational equivalence, with Betti Q-cohomology. A correspondence pi in CH\textasciicircum{}d(X x X)\_Q acts by pi\_*(a)=(p\_2)\_*(p\_1\textasciicircum{}*(a).pi). Murre's full A-D package asserts: (A) there exists a Chow-Kunneth decomposition, namely mutually orthogonal composition-idempotents pi\_0,...,pi\_(2d) in CH\textasciicircum{}d(X x X)\_Q summing to the diagonal and inducing projection onto H\textasciicircum{}i(X,Q) for pi\_i; (B) for EVERY such decomposition and 0<=j<=d, pi\_i acts as zero on CH\textasciicircum{}j when i<j or i>2j; (C) for EVERY such decomposition define F\textasciicircum{}r CH\textasciicircum{}j as the intersection of ker(pi\_i)\_* for max(0,2j-r+1)<=i<=2j, r>=0, with empty intersection F\textasciicircum{}0=CH\textasciicircum{}j; these filtrations are equal for ALL decompositions satisfying A; (D) for EVERY such decomposition, F\textasciicircum{}1 CH\textasciicircum{}j equals ker(CH\textasciicircum{}j(X)\_Q -> H\textasciicircum{}(2j)(X,Q)). Under A and B, F\textasciicircum{}(j+1)=0. No multiplicative or self-dual decomposition, explicit construction algorithm, functorial-extra requirement, or arbitrary-base-field enlargement is included.}
\subsection{Short English statement}
Can algebraic cycles be split into cohomological pieces so that the resulting filtration is independent of the splitting and starts precisely with the homologically trivial cycles?
% BEGIN RESEARCH ADDITION ATTEMPT-153
\subsection{Research attempt}

\paragraph{Scope check and current frontier (24 September 2026).}
This entry uses the \emph{strong} form of B in \sref{153}{1}, Conjecture 5.1:
the forbidden indices include $i<j$, not only $i>2j$.  Its requirement concerning
\emph{every} decomposition is retained.  The full package is known for curves;
for surfaces, constructing a decomposition with B and D must be distinguished
from an assertion about all choices.  Nilpotence and algebraicity of the
cohomological K\"unneth projectors give existence and uniqueness of the motive
summands \emph{up to isomorphism}, not automatically equality of subspaces of
Chow groups (\sref{153}{1}, Proposition 5.3 and Lemma 5.4;
\eref{153}{1}, Theorems 8 and 14).

The recent paper of Laterveer--Vial, published in 2026 from a 2024 manuscript,
records these distinctions and proves cycle-class injectivity on specified
subrings for LLSS eightfolds; it does not prove the full package for arbitrary
varieties (\eref{153}{2}, \S1.1 and Theorem 2).
A relevant established reduction is Vial's theorem that D for products of
curves over a universal domain implies A--D for motives of abelian type
(\eref{153}{3}, Theorem 3).  It would be circular to invoke D for all such
products without an independent proof.

\paragraph{Attack map.}
A lifting approach separates A into algebraicity of the cohomological
projectors and correction of their idempotence and orthogonality in the Chow
ring of correspondences.  The latter is accessible under nilpotence; the
former is still a hypothesis for a general $X$.  A second approach transports
a splitting through correspondences with products of curves, but needs the
corresponding vanishing and homological-kernel assertions on those products.
A third approach attacks the universal quantifier directly: conjugate one
splitting by a homologically trivial perturbation and test whether its
filtration changes.

We pursue the third route in both directions.  It gives a proof strategy if
all relevant perturbations preserve the filtration, and an explicit recipe
for a counterexample if one geometrically realizable perturbation moves it.
The aim is to identify the exact extra condition beyond nilpotence, rather
than mistake isomorphism of motive summands for equality of Chow subgroups.

\paragraph{Attempt I: an explicit conjugacy formula.}
Fix $X$ of dimension $d$ and write
\[
 R_X=\mathrm{CH}^d(X\times X)_{\Q},\qquad
 I_X=\ker\bigl(R_X\longrightarrow H^{2d}(X\times X,\Q)\bigr).
\]
Multiplication in $R_X$ means composition, its unit is $\Delta_X$, and $I_X$
is a two-sided ideal.  K\"unneth and Poincar\'e duality identify the cohomology
class of a degree-zero correspondence with its actions on all $H^i(X,\Q)$.
Thus two Chow--K\"unneth systems $p=(p_i)$ and $q=(q_i)$ satisfy
$q_i-p_i\in I_X$.

For now make the following explicit assumption:
\begin{equation}\label{ra153:nil}
 \text{Every $n\in I_X$ is nilpotent under composition.}
\end{equation}
No exponent uniform in $n$ is required here.  This hypothesis holds when the
Chow motive is Kimura finite-dimensional, by the nilpotence theorem recorded
in \eref{153}{1}, Theorem 8.  The universal finite-dimensionality assertion of
Section~\ref{191} is not being assumed as a known theorem.

Set
\begin{equation}\label{ra153:conjugacy}
 u=\sum_{i=0}^{2d}q_i p_i=1+n,\qquad n\in I_X.
\end{equation}
If $n^N=0$, then $u^{-1}=\sum_{a=0}^{N-1}(-n)^a$.  Orthogonality, with the
order of factors retained, gives
\[
 q_i u=q_i p_i=u p_i,
 \qquad q_i=u p_i u^{-1}.
\]
This is an explicit version of the conjugacy mechanism in
\sref{153}{1}, Lemma 5.4.  It uses only nilpotence of the particular $n$ in
\eqref{ra153:conjugacy}; it does not assume that every sufficiently long
product of different members of $I_X$ vanishes.

Let $V_j=\mathrm{CH}^j(X)_{\Q}$ and let $K_j$ be its homologically trivial
subgroup.  Correspondences act on $V_j$; subscripts $*$ are suppressed below.
For the filtration defined in the original statement, conjugacy gives the
exact identity
\begin{equation}\label{ra153:transport}
 F_q^r V_j=u(F_p^r V_j).
\end{equation}
Indeed, $q_i v=0$ if and only if $p_i u^{-1}v=0$, and the same equivalence
holds for the entire intersection of kernels, including an empty one.

Suppose the reference system $p$ satisfies B and D on $V_j$.  B transfers to
$q$, because a conjugate of the zero operator is zero.  D transfers too:
$I_X V_j\subseteq K_j$ by compatibility of cycle classes with correspondences,
so $u$ and $u^{-1}$ preserve $K_j$ and
$F_q^1V_j=u(K_j)=K_j$.  Thus under \eqref{ra153:nil}, a single verified B,D
splitting gives B,D for \emph{all} splittings on that Chow group.  The
remaining problem is C, not transport of these two properties.

\paragraph{Attempt II: the precise obstruction to C.}
Still assuming \eqref{ra153:nil} and B,D for $p$, the following conditions are
equivalent for the fixed $j$:
\begin{equation}\label{ra153:criterion}
 \begin{aligned}
 &F_q^r V_j=F_p^r V_j\quad\text{for every Chow--K\"unneth system $q$ and $r\geq0$};\\
 &I_X(F_p^r V_j)\subseteq F_p^r V_j\quad\text{for every $r\geq0$}.
 \end{aligned}
\end{equation}
For sufficiency, $u$ in \eqref{ra153:conjugacy} and its finite-series inverse
preserve every $F_p^r$, and \eqref{ra153:transport} gives equality.  Conversely,
for each $n\in I_X$, conjugation by $1+n$ produces another valid system.
Equality of its filtration with $F_p^r$ implies
$(1+n)F_p^r=F_p^r$, hence $nF_p^r\subseteq F_p^r$.  This checks the universal
quantifier, rather than only a preferred class of splittings.

B supplies a finite direct-sum description, valid for every $r\geq0$:
\begin{equation}\label{ra153:split-filtration}
 V_j=\bigoplus_{i=j}^{2j}p_iV_j,\qquad
 F_p^rV_j=\bigoplus_{j\leq i\leq 2j-r}p_iV_j,
\end{equation}
where an empty sum is zero.  Moreover, D and $I_XV_j\subseteq K_j$ imply
$p_{2j}I_XV_j=0$.  Therefore the second line of
\eqref{ra153:criterion} is equivalent to the following finite collection of
\emph{types} of vanishing conditions:
\begin{equation}\label{ra153:upward}
 p_b\,I_X\,p_i(V_j)=0
 \qquad\text{whenever}\qquad j\leq i<b<2j.
\end{equation}
Each condition still quantifies over all correspondences in $I_X$ and all
cycles; this is not a finite computability assertion.
To prove necessity, use the step $r=2j-i$, which contains $p_iV_j$ but no
$p_bV_j$ with $b>i$.  For sufficiency, expand any $n\in I_X$ as
$\sum_{b,i}p_b n p_i$.  The terms with $b>i$ vanish on $V_j$ by
\eqref{ra153:upward}, D, and B; every remaining term preserves all partial sums
in \eqref{ra153:split-filtration}.

In codimension $j=0$ or $1$ there is no pair in
\eqref{ra153:upward}, so no further condition is needed.  In codimension two,
for $d\geq2$, the entire additional condition is
\begin{equation}\label{ra153:codim-two}
 p_3 I_X p_2\bigl(\mathrm{CH}^2(X)_{\Q}\bigr)=0.
\end{equation}
Thus a concrete remaining task is to rule out a correspondence moving a cycle
from the $p_2$ part into the $p_3$ part.  Nilpotence alone has no reason to
forbid such a movement.

\paragraph{Attempt III: an explicit counterexample certificate, if realizable.}
There is an especially small perturbation that would falsify C.  Suppose
$j\leq i<b<2j$ and choose $\gamma\in I_X$ with
\[
 h=p_b\gamma p_i,
 \qquad h(z)\ne0\text{ for some }z\in p_iV_j.
\]
Then $h^2=0$, because $p_i p_b=0$.  No general nilpotence assumption is
needed for this particular construction.  Conjugating by $1+h$ gives
\begin{equation}\label{ra153:elementary-perturbation}
 q_i=p_i+h,\qquad q_b=p_b-h,\qquad q_a=p_a\ (a\notin\{i,b\}).
\end{equation}
Direct multiplication verifies orthogonality and idempotence.  Their sum and
cohomological actions are unchanged.  At $r=2j-i$, the cycle $z$ belongs to
$F_p^r$, but $q_bz=-h(z)\ne0$, so $z\notin F_q^r$.  This would be an actual
counterexample to the printed universal C, not merely nonuniqueness of its
individual projectors.  No such geometric $h,z$ has been constructed here.

A small algebraic model shows that the obstacle is real at the ring level.
Take the ring and its nilpotent ideal
\[
 R=\left\{\begin{pmatrix}a&0&0\\ b&c&0\\0&0&d\end{pmatrix}:a,b,c,d\in\Q\right\},
 \qquad I=\Q E_{21},\qquad I^2=0,
\]
acting on $V=\Q e_2\oplus\Q e_3\oplus\Q e_4$, where the basis subscripts are
weight labels, not matrix indices.  Put $p_2=E_{11}$, $p_3=E_{22}$,
$p_4=E_{33}$ and $p_0=p_1=0$.  Give the degree-four cycle-class model the map
$\mathrm{cl}(x e_2+y e_3+z e_4)=z$.  Then
\[
 F_p^1=\ker(\mathrm{cl})=\Q e_2\oplus\Q e_3,
 \qquad F_p^2=\Q e_2,
 \qquad F_p^3=0.
\]
For $h=E_{21}$, \eqref{ra153:elementary-perturbation} gives
\[
 q_2=\begin{pmatrix}1&0&0\\1&0&0\\0&0&0\end{pmatrix},\quad
 q_3=\begin{pmatrix}0&0&0\\-1&1&0\\0&0&0\end{pmatrix},\quad q_4=p_4,
 \qquad F_q^2=\Q(e_2+e_3)\ne F_p^2.
\]
Here $F_q^1=F_p^1$, so the analogue of D survives, as does B.  Exact rational
matrix multiplication verifies the displayed identities; the failure of C
is also immediate by applying $q_3$ to $e_2$.
This is \emph{not} a smooth projective variety or a counterexample to Murre.
It specifically refutes the proposed purely ring-theoretic step
``nilpotence plus B,D makes the filtration independent.''

\paragraph{Attempt IV: searching for a geometric perturbation.}
The simplest construction to test is a sum of decomposable correspondences
$\alpha\times\beta$, with
$\alpha\in\mathrm{CH}^{d-s}(X)_{\Q}$ and
$\beta\in\mathrm{CH}^{s}(X)_{\Q}$.  On $V_j$ it acts as zero unless $s=j$,
and for $s=j$ its action is
\begin{equation}\label{ra153:rank-one}
 (\alpha\times\beta)_*(z)=\deg(z\cdot\alpha)\,\beta.
\end{equation}
This follows directly by integrating the first factor in the pullback,
intersection, and pushforward formula.  If $z\in p_iV_j$ with $i<2j$, its
cycle class is zero: the cycle class lies in $H^{2j}$ and $p_i$ projects to
$H^i$.  Hence $\deg(z\cdot\alpha)=0$ for every such $\alpha$.  Thus every
candidate built only from these external products kills the input in
\eqref{ra153:upward}.  Choosing $\beta$ homologically trivial does not help.
This construction fails by an exact intersection-number calculation, not by
lack of a numerical example.  A successful negative attack would need a
non-decomposable correspondence acting nontrivially on homologically trivial
cycles in the indicated direction.

For the positive attack, the known product argument explains why the missing
vanishing is deep.  Form the Chow--K\"unneth system on $X\times X$
\[
 P_m=\sum_{a+c=m}{}^t p_{2d-a}\times p_c,
 \qquad 0\leq m\leq4d,
\]
with the factors rearranged in the usual way to make a correspondence on
$X\times X$.  Pull--push composition gives
\[
 ({}^t p_i\times p_b)_*(\gamma)=p_b\gamma p_i.
\]
Its total cohomological index is $2d-i+b$.  B in codimension $d$ on
$X\times X$, for this system, would kill it when $b>i$.  D there would
additionally kill $p_i I_X p_i$.  Consequently $I_X$ would be strictly
lower triangular, and $I_X^{2d+1}=0$ by counting successive decreases of
indices.  This is the mechanism of \sref{153}{1}, Proposition 5.8, not a new
proof of its hypotheses.  It would make every condition
\eqref{ra153:upward} automatic, but assumes the requisite statements on the
square rather than deriving them from statements on $X$.

\paragraph{Stress test and dependency check.}
All equalities in the perturbation calculation take place modulo rational
equivalence, except where explicitly passed to cohomology.  The matrix model
is used only to falsify a proposed abstract lemma; no geometric realization,
Poincar\'e duality, or Chow-motive interpretation of that model is asserted.
The calculations use finite sums and finite nilpotent inverses; there is no
convergence or unproved completion of the correspondence ring.  The perturbed
systems need not be self-dual or multiplicative; neither extra condition is
part of this entry.

The proof of the criterion handles every admissible $j$, every $r\geq0$,
and every Chow--K\"unneth system.  In particular, the indices $i<j$ in the
strong B have not been omitted.  Its assumptions still include one actual
splitting with B,D.  They do not establish A for arbitrary $X$, nor prove
that the first splitting has B,D.

Assuming Kimura finite-dimensionality for all $X$ would use the unresolved
Section~\ref{191}; nilpotence here concerns \emph{composition}, not the
external powers in Section~\ref{182}.  Algebraizing every K\"unneth projector
would itself establish the K\"unneth standard conjecture, not just a formal
linear-algebra step.  Finally, using a full Bloch--Beilinson filtration to
justify invariance would invoke a conjectural structure equivalent to the
universal Murre package in the framework of \sref{153}{1}, Theorem 5.2.
None of these dependencies is discharged by the calculation above.

\paragraph{Outcome.}
\textbf{Outcome: Conditional reduction.}

For a fixed $X$, suppose A holds, every element of $I_X$ is composition-nilpotent,
and one chosen splitting satisfies strong B and D for every codimension.
Then the \emph{full all-decompositions} A--D package is equivalent to the
vanishings \eqref{ra153:upward} for that splitting.  The implication in each
direction has been proved above; no priority claim is made for the underlying
conjugacy and weight-filtration methods.  In codimension two, the additional
condition is exactly \eqref{ra153:codim-two}.  A nonzero geometric instance of
that upward action would instead give the explicit counterexample certificate
\eqref{ra153:elementary-perturbation}.

Nilpotence alone is insufficient, as the matrix model shows, and the most
obvious external-product realization of the bad perturbation is impossible
by \eqref{ra153:rank-one}.  What remains is the genuinely geometric vanishing
or construction of such actions, together with the separate existence and
B,D hypotheses.  No smooth projective counterexample or unconditional proof
of the universal target has been obtained.
% END RESEARCH ADDITION ATTEMPT-153
\subsection{Sources}
\begin{itemize}\raggedright
\item[\textnormal{[S1]}]\hypertarget{source:153:1}{} Uwe Jannsen, Motivic Sheaves and Filtrations on Chow Groups, Proc.Symp.Pure Math.55(1994),Part1,245-302.
\url{https://epub.uni-regensburg.de/26639/1/ubr13239_ocr.pdf}.
{\small\textit{Catalogue formulation source.}}
% BEGIN RESEARCH ADDITION REFERENCES-153
\item[\textnormal{[E1]}]\hypertarget{extra:153:1}{} Vladimir Guletskii and
Claudio Pedrini, \emph{Finite dimensional motives and the conjectures of
Beilinson and Murre}, K-Theory 30(3) (2003), 243--263.
Theorems 8 (Kimura's nilpotence theorem) and 14 (isomorphism of the summands)
in the consulted manuscript:
\url{https://arxiv.org/html/math/0303170}.
{\small\textit{Consulted 24 September 2026; isomorphism of motives is not
identified with equality of their images in a fixed Chow group.}}
\item[\textnormal{[E2]}]\hypertarget{extra:153:2}{} Robert Laterveer and Charles
Vial, \emph{The Beauville--Voisin--Franchetta conjecture and LLSS eightfolds},
Indagationes Mathematicae 37 (2026), 250--270;
arXiv:2404.10465v1 (16 April 2024), \S1.1, Theorem 2, Proposition 1.7 and
Corollary 1.10.
\url{https://arxiv.org/html/2404.10465v1}.
Publication information:
\url{https://www.math.uni-bielefeld.de/~vial/}.
{\small\textit{Manuscript and author's publication list consulted
24 September 2026. The result on selected subrings is not used as an
injectivity theorem for the entire Chow ring.}}
\item[\textnormal{[E3]}]\hypertarget{extra:153:3}{} Charles Vial,
\emph{Remarks on motives of abelian type}, arXiv:1112.1080v3
(27 July 2015), Theorem 3 and \S2; published in Tohoku Mathematical Journal
69(2) (2017).
\url{https://arxiv.org/html/1112.1080v3}.
{\small\textit{Consulted 24 September 2026. The assumption of Murre (D)
for products of curves over a universal domain is retained.}}
% END RESEARCH ADDITION REFERENCES-153
\end{itemize}

\end{document}
